\section{The FIBER compiler and its main theorem}
\label{sec:compiler}

\subsection{A small proof-producing intermediate language}

The source language is not intended to replace Lean or the mathematical prose
of a packet.  It records only the finite-use structure needed for compilation.
A program is a finite directed graph whose nodes have one of the following
forms.

\begin{longtable}{p{0.18\linewidth}p{0.72\linewidth}}
\toprule
Node & Meaning\\
\midrule
\endfirsthead
\toprule
Node & Meaning\\
\midrule
\endhead
\textsf{primitive} & Apply a fixed primitive-recursive operation and a local checked contract.\\
\textsf{pair}/\textsf{project} & Form or read a finite tuple or encoded sequence.\\
\textsf{map}/\textsf{fold} & Execute a bounded loop with a quantifier-free state invariant.\\
\textsf{search} & Use checked bounded least search for a reflected verifier.\\
\textsf{case} & Branch on a reflected Boolean atom and join the two derivations by disjunction elimination.\\
\textsf{exists} & Open a local certificate and seal it by existential elimination.\\
\textsf{strategy} & Verify and realize a fixed bounded quantifier game.\\
\textsf{observe} & Export an equivalence-invariant or metric-stable observation.\\
\textsf{residual} & Add an exact producer formula to the local context.\\
\textsf{call} & Invoke another compiled interface at a finite list of demanded queries.\\
\bottomrule
\end{longtable}

Each node declares its input and output arities, query transformer, certificate
bound, verifier bit, formula template, and local proof theorem.  Acyclic graphs
are topologically scheduled.  A graph with monotone cyclic demand first invokes
\cref{thm:demand-closure}; the resulting finite demand is then unrolled into an
acyclic schedule.

\subsection{The compiler invariant}

For every scheduled node \(v\), the compiler maintains:

\begin{enumerate}
\item a term tuple \(s_v\) containing the node state;
\item a quantifier-free state formula \(I_v(s_v)\);
\item a finite list \(\Gamma_v\) of residual and input formulas;
\item a derivation \(\Gamma_v\vdash_{\PRAAx}I_v(s_v)\);
\item a map from local certificate variables to their lexical scopes;
\item an observation relation guaranteeing that no noninvariant local code
      escapes.
\end{enumerate}

Bounded logical content is reflected into bits before entering a loop.
Quantifiers may occur in the external contract of a node, but the formula on
which the \(\ipra\) induction rule acts is always quantifier-free.

\subsection{Compilation theorem}

\begin{theorem}[FIBER compilation]
\label{thm:compilation}
Let \(P\) be a fixed well-formed FIBER program from interfaces
\(X_1,\ldots,X_m\) to \(Y\).  Suppose:

\begin{enumerate}
\item every primitive node is supplied with a local derivation in
      \(\Prf\;\PRAAx\);
\item every bounded formula used by a checker has a reflection derivation;
\item every loop carries a quantifier-free invariant with base and step
      derivations;
\item every cyclic demand operator is monotone and has explicit finite caps;
\item every exported operation respects the declared observation relation;
\item every external requirement is named by a residual formula.
\end{enumerate}

Then an effective metalevel compiler constructs:

\begin{enumerate}
\item fixed primitive-recursive symbols for demand, master-certificate bounds,
      verification, transformation, and observation;
\item a finite list of residual formulas \(R_P(p,q)\);
\item a Lean proof term of the literal type
\[
 R_P(p,q),\;\mathsf{Inputs}_P(p,q)
 \;\vdash_{\PRAAx}\;
 \mathsf{Output}_Y(p,q);
\]
\item and, when the output is certificate-independent, a proof in which all
      local and master certificates have been existentially eliminated.
\end{enumerate}

If derivations of all formulas in \(R_P(p,q)\) are later supplied, the closed
or client-context conclusion follows by \texttt{Prf.trans}.
\end{theorem}

\begin{proof}
The proof is by external induction over the finite schedule.

At a primitive node, instantiate its closed local theorem at the actual input
terms using \texttt{LeanProofs.IPRA.inst}, then weaken to the current context.
At a pairing or projection node, use the checked pairing equations.  At a
bounded map or fold, apply quantifier-free induction to the declared state
formula; the base and step derivations are part of the node contract.  At a
search node, instantiate \cref{thm:least}.  At a case node, reflect the bit,
derive the finite decidable disjunction, compile each branch, and join them by
\texttt{Prf.orE}.  At an existential node, compile the body under the shifted
certificate hypothesis and apply \texttt{Prf.exE}.  At a strategy node, use
\cref{thm:strategy}.  At an observation node, apply the operation congruence
or metric stability theorem and then seal the local certificate.

A call node first computes the finite demanded queries.  For each demand it
instantiates the called interface theorem.  The input derivations are then
substituted into the caller's local theorem using \texttt{Prf.trans}.  If a
single code is required, \cref{thm:master} forms the master certificate; if
not, the certificates are opened one at a time and sealed.

For monotone cyclic demand, \cref{thm:demand-closure} computes the least finite
closure.  The graph is unrolled for the resulting bounded number of instances,
so the remaining schedule is finite and acyclic.

The induction maintains the lexical-scope invariant.  A certificate variable
can disappear only through existential elimination; an exported term may
mention it only if the node is marked canonical and least search has produced
a closed term.  Therefore the final conclusion contains no accidental witness
escape.  Every unresolved producer encountered by the schedule is inserted as
a hypothesis, yielding exactly \(R_P(p,q)\).  Finally simultaneous cut
combines all node derivations and, when available, discharges residuals.
\end{proof}

\subsection{Composition}

\begin{theorem}[Compositionality]
\label{thm:composition}
If \(R:X\to Y\) and \(S:Y\to Z\) are FIBER realizers, then there is a FIBER
realizer \(S\circ R:X\to Z\).  Its demand is obtained by flattening the
\(R\)-demands for each \(Y\)-query demanded by \(S\); its residual list is the
union of the specialized residuals of \(R\) and \(S\); and its soundness proof
is obtained by substitution and simultaneous cut.
\end{theorem}

\begin{proof}
For a target query \(q\), compute \(\dem_S(q)=\langle y_0,\ldots,y_{k-1}\rangle\).
For each \(y_i\), compute \(\dem_R(y_i)\), concatenate the lists, and remove
no duplicates unless a proved normalization theorem is available.  A master
certificate supplies the required \(X\)-certificates.  The transformer first
runs \(T_R\) at each \(y_i\), bundles the resulting \(Y\)-certificates, and
then runs \(T_S\).  Bounds and all operations are primitive recursive.

In the object theory, instantiate the soundness theorem for \(R\) at every
\(y_i\).  Use these derivations to discharge the \(Y\)-validity hypotheses of
the soundness theorem for \(S\) by \texttt{Prf.trans}.  Extensionality or
metric stability composes by transitivity or the triangle inequality.  Local
certificates are sealed.  The remaining hypotheses are precisely the
specialized residuals.  Sequence concatenation may make associativity hold
only up to a provable coding equality, but the exported observation is
unchanged; normalizing all flattened demands to one canonical sequence code
makes composition strictly associative at the metadata level.
\end{proof}

\begin{corollary}[Reuse under arbitrary terms and contexts]
A compiled closed interface theorem can be specialized at arbitrary object
terms and used under arbitrary client hypotheses.
\end{corollary}

\begin{proof}
Apply \texttt{LeanProofs.IPRA.inst} to substitute the terms and then
\texttt{Prf.ofNil} or weakening to add the client context.  This is why reusable
lemmas should be stated either at arbitrary terms or in an explicitly
universally closed form.
\end{proof}

\subsection{Soundness versus completeness}

The compilation theorem is a soundness theorem: every accepted program yields
an \(\ipra\) derivation from its explicit residuals.  It does not assert that
every \(\ipra\) theorem has a convenient FIBER presentation, nor that every
classical proof can be mined automatically.  This asymmetry is intentional.
The source language is a disciplined fragment optimized for finite,
quantitative mathematics and auditability.

\subsection{Why the framework is not merely a checker generator}

A checker alone proves only a conditional statement of the form ``if this
finite code verifies, then the local conclusion holds.''  The compiler also
tracks:

\begin{itemize}
\item where the certificate comes from;
\item whether its existence is local, canonical, or residual;
\item how a downstream query changes upstream precision demand;
\item whether two independently produced certificates are compatible;
\item which observations are stable under representation changes;
\item which formulas are actually used as quantifier-free induction
      invariants;
\item and how all local proofs are assembled into the final object sequent.
\end{itemize}

These are exactly the points at which the existing packet outlines most often
stop.
